\chapter{Error messages}
\label{ch:diagnostics}
\index{error messages}\index{warnings}

Every message is printed on standard error. An \texttt{error:} stops the
program with exit status~1 (output printed so far is kept); a
\texttt{warning:} does not. Messages for a file name the file, and for
lexical and syntax errors the line, as \textit{file}\texttt{:}\textit{line}.
Both back ends report the same errors; the native compiler reports static
errors at compile time and run-time errors when the executable runs.

\section{Lexical errors}

\begin{center}\small
\begin{tabularx}{\textwidth}{@{}l X@{}}
\toprule
\textbf{Message} & \textbf{Cause} \\
\midrule
\tc{unexpected character 'c'} & a character that begins no token, such as \tc{@} or \tc{&} alone \\
\tc{unterminated string literal} & a string without its closing quote \\
\tc{malformed character literal} & \tc{'ab'}, \tc{''} or an unknown escape \\
\tc{unterminated block comment} & a \tc{/*} without \tc{*/} \\
\tc{integer literal n is too large} & a literal that does not fit (\S\ref{sec:int-width}) \\
\bottomrule
\end{tabularx}
\end{center}

\forgeexample*{error-lex}

\section{Syntax errors}

\begin{center}\small
\begin{tabularx}{\textwidth}{@{}l X@{}}
\toprule
\textbf{Message} & \textbf{Cause} \\
\midrule
\tc{syntax error near `t`} & the parser could not continue at token \tc{t}; the line is that of the token \\
\tc{unexpected end of input} & the file ends in the middle of a definition \\
\bottomrule
\end{tabularx}
\end{center}

\forgeexample*{error-syntax}

\section{Imports, modules, traits and macros}

\begin{center}\small
\begin{tabularx}{\textwidth}{@{}l X@{}}
\toprule
\textbf{Message} & \textbf{Cause} \\
\midrule
\tc{import not found: f} & no file \tc{f} next to the importing file \\
\tc{import cycle detected at f} & files that import each other \\
\tc{duplicate inter/trait: T} & a trait declared twice \\
\tc{impl missing method: m} & an \tc{impl} without method \tc{m} (and \tc{m} has no default) \\
\tc{impl has unknown method: m} & an \tc{impl} defining a method that its trait does not declare \\
\tc{duplicate method in impl} & a method defined twice in one \tc{impl} \\
\tc{duplicate impl of T for X} & two \tc{impl}s of a trait for the same type \\
\tc{unknown macro m!} & invocation of a macro that is not defined \\
\tc{no matching arm for macro m!} & no rule of the macro accepts the arguments \\
\bottomrule
\end{tabularx}
\end{center}

\forgeexample*{error-impl}

\section{Type errors}

Type errors start with \texttt{type error:}. A mismatch names the two types
that could not be made equal and, on the next line, the definition (as the
type checker sees it, with operators as applications) in which it was found.

\begin{center}\small
\begin{tabularx}{\textwidth}{@{}l X@{}}
\toprule
\textbf{Message} & \textbf{Cause} \\
\midrule
\tc{Unbound variable: x} & \tc{x} is not defined (or not yet: definitions are in order) \\
\tc{Unknown type: T} & an annotation names an undeclared type \\
\tc{Unknown constructor in pattern: C} & a pattern uses an undeclared constructor \\
\tc{... cannot match A with B} & a value of type \tc{B} is used where \tc{A} is needed \\
\tc{No instance `C` for type T} & a trait method or constrained function used at a type without an \tc{impl} \\
\tc{the variable `x` is bound twice in the same pattern} & a pattern that binds a name twice \\
\bottomrule
\end{tabularx}
\end{center}

\paragraph{Ambiguous types} When a trait method is used at a type that
nothing in the program determines --- \fc|show empty|, where \fc|empty|
could be the empty string, list or option --- no instance can be chosen.
The interpreter reports \texttt{runtime error: ambiguous type: cannot tell
which `Show` instance `show` should use here} when the method is called;
the native compiler rejects the program (\texttt{the type at which `Monoid`
is used is ambiguous here}). An annotation resolves it:
\fc|let e : String = empty|, then \fc|show e|.

\forgeexample*{error-type}

\forgeexample*{error-unbound}

\forgeexample*{error-ctor}

\forgeexample*{error-instance}

\section{Run-time errors}

Run-time errors start with \texttt{runtime error:}.

\begin{center}\small
\begin{tabularx}{\textwidth}{@{}l X@{}}
\toprule
\textbf{Message} & \textbf{Cause} \\
\midrule
\tc{division by zero} & \tc{/} or \tc{\%} with a zero divisor \\
\tc{no pattern matched} & no arm of a \tc{case}, or a refutable \tc{let} or parameter pattern, matched; \tc{list_head}, \tc{list_tail} and \tc{list_nth} outside the list \\
\tc{cannot compare functions} & \tc{==} or \tc{<} on functions \\
\tc{stack overflow} & recursion too deep (interpreter) \\
\bottomrule
\end{tabularx}
\end{center}

\forgeexample{error-division}

\section{Warnings}

\begin{center}\small
\begin{tabularx}{\textwidth}{@{}l X@{}}
\toprule
\textbf{Message} & \textbf{Cause} \\
\midrule
\tc{in `f`: this case is not exhaustive: `v` is not matched} & a \tc{case} in definition \tc{f} has no arm for value \tc{v} \\
\tc{in `f`: redundant case arm: `p` can never match} & an arm that the arms before it cover \\
\bottomrule
\end{tabularx}
\end{center}

\noindent Section~\ref{sec:exhaustiveness} shows both.
